% Current published PDF pp33:2--3,33:5--6: exact syntax and derivation context.
\section*{Higher-order infinitary syntax and deduction}
\begin{authordef}[1]
Types and terms are defined inductively:
\[\begin{array}{rcl}\tau,\sigma&::=&\iota\mid o\mid o^I\mid\tau\to\sigma,\\t,u,A_i,F&::=&x\mid c\mid t\,u\mid\lambda x.t\mid\{i\mapsto A_i\}_I\mid\pi_jF.\end{array}\]
Here $\iota$ is the type of individuals, $o$ is the type of propositions, $o^I$ is the type of families of propositions indexed by a countable set $I$, $x$ is a variable, $c$ is a constant, $\{i\mapsto A_i\}_I$ is the family of propositions $(A_i)_{i\in I}$, and $\pi_jF$ is the projection on the $j$-th component of $F$.
\end{authordef}
The typing rules are the usual ones, extended with: if for every $i\in I$ we have $A_i$ of type $o$ then $\{i\mapsto A_i\}_I$ is of type $o^I$; if $F$ is of type $o^I$ then $\pi_jF$ is of type $o$. We write $\Lambda_\tau$ for the set of terms of type $\tau$. We restrict terms so that they do not have infinitely many free variables. We write $t[x\leftarrow u]$ for the term $t$ in which $x$ has been substituted by $u$.
\begin{authordef}[3]
We consider the following rules:
\[(\lambda x.t)u\hookrightarrow t[x\leftarrow u],\quad\pi_j\{i\mapsto A_i\}_I\hookrightarrow A_j,\quad\lambda x.tx\hookrightarrow t\ (x\notin\mathrm{fv}(t)),\quad\{i\mapsto\pi_iF\}_I\hookrightarrow F.\]
The reduction relation $\hookrightarrow$ is the smallest relation that contains these rules and such that:
\begin{itemize}\item if $t\hookrightarrow t'$, then $tu\hookrightarrow t'u$, $ut\hookrightarrow ut'$ and $\lambda x.t\hookrightarrow\lambda x.t'$;
\item if $A_i\hookrightarrow B_i$ for every $i\in J$, where $J\ne\varnothing$ and $J\subseteq I$, and $A_i=B_i$ otherwise, then $\{i\mapsto A_i\}_I\hookrightarrow\{i\mapsto B_i\}_I$;
\item if $F\hookrightarrow F'$ then $\pi_jF\hookrightarrow\pi_jF'$.
\end{itemize}
The relation $\equiv$ is the reflexive, symmetric and transitive closure of $\hookrightarrow$. The $\beta\pi$-reduction $\hookrightarrow_{\beta\pi}$ and $\beta\pi$-conversion $\equiv_{\beta\pi}$ are defined like $\hookrightarrow$ and $\equiv$, but only consider the $\beta$-reduction rule and the projection rules.
\end{authordef}
There are particular constants, called logical constants, for defining the connectives and quantifiers: contradiction $\bot$ of type $o$, implication $\Rightarrow$ of type $o\to o\to o$, quantifiers $\forall_\tau$ and $\exists_\tau$ of type $(\tau\to o)\to o$, infinite conjunction $\bigwedge_I$ and infinite disjunction $\bigvee_I$ of type $o^I\to o$. Terms $\forall_\tau(\lambda x.A)$ and $\exists_\tau(\lambda x.A)$ are abbreviated as $\forall x.A$ and $\exists x.A$. Negation is $\neg A:=A\Rightarrow\bot$ and $A\Leftrightarrow B:=(A\Rightarrow B)\wedge(B\Rightarrow A)$. Formulas are terms of type $o$.
A context $\Gamma$ is a finite sequence of formulas. We write $\Gamma\vdash A$ when $A$ is derivable from $\Gamma$. We write $\mathrm{fv}(t_1,\ldots,t_n)$ for the set of free variables that occur in $t_1,\ldots,t_n$. The ordinary natural-deduction rules are obtained by ignoring the cut-free/neutral adornments of Figure1 in the original paper. The finite higher-order natural-deduction rules, conversions and monad facts are retained in the complete collected author precursor Sections2--6 below. The additional infinitary rules, with those adornments erased as specified by the authors, are:
\[\frac{\Gamma\vdash\pi_iF\text{ for every }i\in I}{\Gamma\vdash\bigwedge_I F}\ \mathrm{AndInf\!\!\!-I}\qquad\frac{\Gamma\vdash\bigwedge_I F}{\Gamma\vdash\pi_jF}\ \mathrm{AndInf\!\!\!-E}\]
\[\frac{\Gamma\vdash\pi_jF}{\Gamma\vdash\bigvee_I F}\ \mathrm{OrInf\!\!\!-I}\qquad\frac{\Gamma\vdash\bigvee_I F\quad\Gamma,\pi_iF\vdash C\text{ for every }i\in I}{\Gamma\vdash C}\ \mathrm{OrInf\!\!\!-E}.\]
There are versions of AndInf-E and OrInf-I for every $j\in I$. Derivation trees possibly have an infinite width, but each branch has a finite height. The formula $A\wedge B$ (respectively $A\vee B$) is defined by $\bigwedge_I\{i\mapsto A_i\}_{\{0,1\}}$ (respectively $\bigvee_I\{i\mapsto A_i\}_{\{0,1\}}$), where $A_0=A$ and $A_1=B$. The finite conjunction and disjunction rules are particular cases of those for the infinite conjunction and disjunction. We consider either classical logic $\CL$ or intuitionistic logic $\IL$ -- in which we do not have the PEM rule. Weakening, contraction and exchange are admissible.
% Current original Section5 pp33:14--17.
\section*{Barr's theorem for higher-order geometric theories}
Geometric formulas are built using finite conjunctions, infinite disjunctions and existential quantifiers. Geometric implications are formulas of the form $\forall\vec x.(A\Rightarrow B)$, where $A$ and $B$ are geometric formulas and $\vec x$ is a finite sequence of variables (possibly empty). A theory $\mathcal T$ is geometric when all its axioms are geometric implications. Throughout this section, we consider a higher-order infinitary logic with infinite disjunctions, but with finite conjunctions.
\begin{authordef}[39, Friedman's translation]
Let $A$ be a formula in higher-order infinitary logic. Its Friedman's translation is $A^T:=TA_T$, where $t\mapsto t_T$ is inductively defined by:
\[\begin{array}{rcl}x_T&:=&x,\\c_T&:=&\begin{cases}\lambda p.\forall x.T(px)&\text{if }c=\forall,\\\lambda p.\lambda q.p\Rightarrow Tq&\text{if }c=\Rightarrow,\\c&\text{otherwise},\end{cases}\\(\lambda x.t)_T&:=&\lambda x.t_T,\\(tu)_T&:=&t_Tu_T,\\(\{i\mapsto A_i\}_I)_T&:=&\{i\mapsto(A_i)_T\}_I,\\(\pi_jF)_T&:=&\pi_jF_T.\end{array}\]
\end{authordef}
As we are in higher-order logic, our translation must commute with substitution and preserve convertibility. Note that all the results of this section also work if we only consider the $\beta$-reduction rule and the projection rules.
\begin{authorprop}[40]
Let $t,u$ be terms and $A,B$ be higher-order infinitary formulas.
\begin{enumerate}\item $(t[z\leftarrow w])_T=t_T[z\leftarrow w_T]$.\item $(A[z\leftarrow w])^T=A^T[z\leftarrow w_T]$.\item If $t\equiv u$ then $t_T\equiv u_T$.\item If $A\equiv B$ then $A^T\equiv B^T$.\end{enumerate}
\end{authorprop}
\begin{proof}The first item is proved by induction on the term $t$. The second item is a corollary of the first item. The third item is proved by induction on $t\equiv u$. The fourth item is a corollary of the third item.\end{proof}
Let us remark that the operator $TA:=(A\Rightarrow R)\Rightarrow R$ corresponds to a monad operator in the sense of original reference38, meaning that it satisfies the (unit) and (bind) properties:
\[\vdash_L A\Rightarrow TA,\tag{unit}\]
\[\vdash_L(A\Rightarrow TB)\Rightarrow(TA\Rightarrow TB).\tag{bind}\]
It follows that the following properties hold [38].
\begin{authorprop}[41]
Let $A,B$ be formulas and $P$ be a predicate.
\begin{enumerate}\item $\vdash_\IL TTA\Leftrightarrow TA$.\item $\vdash_\IL(A\Rightarrow B)\Rightarrow(TA\Rightarrow TB)$.\item $\vdash_\IL T(A\wedge B)\Leftrightarrow(TA\wedge TB)$.\item $\vdash_\IL T(A\Rightarrow B)\Rightarrow(TA\Rightarrow TB)$.\item $\vdash_\IL T(A\Rightarrow TB)\Leftrightarrow(TA\Rightarrow TB)$.\item $\vdash_\IL T(\forall x.T(Px))\Leftrightarrow\forall x.T(Px)$.\item $\vdash_\IL T(\exists x.T(Px))\Leftrightarrow T(\exists P)$.\end{enumerate}
\end{authorprop}
\begin{authortheorem}[42]
Let $TA=(A\Rightarrow R)\Rightarrow R$ with parameter $R$. Let $\Gamma,A$ be higher-order infinitary formulas. Let $\mathcal T$ be a higher-order geometric theory. If $\Gamma\vdash_\CL A$ in $\mathcal T$ then $\Gamma^T\vdash_\IL A^T$ in $\mathcal T$.
\end{authortheorem}
\begin{proof}
We proceed by induction on the derivation. Most of the cases are direct consequences of Proposition41 and are detailed in [38, see Lemma16]. We only show the important cases.
\begin{description}
\item[OrInf-I:] By induction hypothesis, we have $\Gamma^T\vdash_\IL T(\pi_jF)_T$. Using Ax, we have $\Gamma^T,(\pi_jF)_T\vdash_\IL(\pi_jF)_T$. Using OrInf-I, we get $\Gamma^T,(\pi_jF)_T\vdash_\IL\bigvee_I F$. Using Imp-I and Proposition41(2), we derive $\Gamma^T\vdash_\IL T(\pi_jF)_T\Rightarrow T(\bigvee_I F)$. We conclude using the induction hypothesis.
\item[OrInf-E:] By induction hypotheses, we have $\Gamma^T\vdash_\IL T(\bigvee_I F)$ and $\Gamma^T,(\pi_iF)_T\vdash_\IL TC_T$ for every $i\in I$. By Ax, we have $\Gamma^T,\bigvee_I F\vdash_\IL\bigvee_I F$. By weakening of the induction hypotheses, we have $\Gamma^T,\bigvee_I F,(\pi_iF)_T\vdash_\IL TC_T$ for every $i\in I$. Therefore, we derive $\Gamma^T,\bigvee_I F\vdash_\IL TC_T$ using OrInf-E. We get $\Gamma^T\vdash_\IL T(\bigvee_I F)\Rightarrow TC_T$ using Imp-I and (bind). We conclude using Imp-E with the first induction hypothesis.
\item[PEM:] We prove $\Gamma^T\vdash_\IL T(A\vee\neg A)_T$ using the following author derivation (the rows retain the original successive inference labels):
\[\begin{array}{rcl}\mathrm{Ax}:&\Gamma^T,B\Rightarrow R,A_T,\bot\Rightarrow R&\vdash_\IL A_T\\\mathrm{Or\!\!-IL}:&\Gamma^T,B\Rightarrow R,A_T,\bot\Rightarrow R&\vdash_\IL B\\\mathrm{Imp\!\!-E\ with}\ B\Rightarrow R:&\Gamma^T,B\Rightarrow R,A_T,\bot\Rightarrow R&\vdash_\IL R\\\mathrm{Imp\!\!-I}\times2:&\Gamma^T,B\Rightarrow R&\vdash_\IL A_T\Rightarrow((\bot\Rightarrow R)\Rightarrow R)\\\mathrm{Or\!\!-IR}:&\Gamma^T,B\Rightarrow R&\vdash_\IL B\\\mathrm{Imp\!\!-E\ with}\ B\Rightarrow R:&\Gamma^T,B\Rightarrow R&\vdash_\IL R\\\mathrm{Imp\!\!-I}:&\Gamma^T&\vdash_\IL((A\vee(A\Rightarrow((\bot\Rightarrow R)\Rightarrow R)))\Rightarrow R)\Rightarrow R\end{array}\]
where $B$ is an abbreviation for the formula $A_T\vee(A_T\Rightarrow((\bot\Rightarrow R)\Rightarrow R))$.
\end{description}
\end{proof}
\paragraph{Original Remark43.} The And-I case crucially depends on $\vdash_\IL(TA\wedge TB)\Rightarrow T(A\wedge B)$. The proof of this property does not extend to the infinite case. This is why, in this section, we consider finite conjunctions and not infinite conjunctions.
Geometric formulas do not contain implications or universal quantifiers, while the translation $t\mapsto t_T$ only inserts $T$ after universal quantifiers and implications. It follows that for any geometric formula $A$, we directly have $A_T=A$. The same result does not hold for geometric implications $\forall\vec x.(A\Rightarrow B)$. But geometric implications imply their translation.
\begin{authorlemma}[44]
Let $A,B$ be higher-order geometric formulas, and $\vec x$ be a finite sequence of variables (possibly empty). We have $\vdash_\IL(\forall\vec x.(A\Rightarrow B))\Rightarrow(\forall\vec x.(A\Rightarrow B))_T$.
\end{authorlemma}
\begin{proof}
Since $(A\Rightarrow B)_T=A\Rightarrow TB$, we directly have $\vdash_\IL(A\Rightarrow B)\Rightarrow(A\Rightarrow B)_T$ using (unit) on $B$. If $\vec x$ is empty, this achieves the case. Suppose that $\vec x=x_1,\ldots,x_n$. We derive $\vdash_\IL(A\Rightarrow B)\Rightarrow T(A\Rightarrow B)_T$ using (unit) on $(A\Rightarrow B)_T$. Using All-I, we derive $\vdash_\IL\forall x_n.((A\Rightarrow B)\Rightarrow T(A\Rightarrow B)_T)$. Then we use the distributivity of $\forall$ over $\Rightarrow$ to derive $\vdash_\IL(\forall x_n.(A\Rightarrow B))\Rightarrow(\forall x_n.T(A\Rightarrow B)_T)$, that is $\vdash_\IL(\forall x_n.(A\Rightarrow B))\Rightarrow(\forall x_n.(A\Rightarrow B))_T$. We iterate this procedure to get $\vdash_\IL(\forall\vec x.(A\Rightarrow B))\Rightarrow(\forall\vec x.(A\Rightarrow B))_T$.
\end{proof}
% MAIN STATEMENT BEGIN -- current official Theorem45 p33:17.
\begin{authortheorem}[45]
Let $\mathcal T$ be a higher-order geometric theory, $A$ and $B$ be higher-order geometric formulas, and $\vec x$ be a finite sequence of variables (possibly empty). If $\vdash_\CL\forall\vec x.(A\Rightarrow B)$ in $\mathcal T$ then $\vdash_\IL\forall\vec x.(A\Rightarrow B)$ in $\mathcal T$.
\end{authortheorem}
% MAIN PROOF BEGIN.
\begin{proof}
Suppose that $\vdash_\CL\forall\vec x.(A\Rightarrow B)$ in $\mathcal T$. Then there exists some geometric implications $\Theta$ of $\mathcal T$ such that $\Theta\vdash_\CL\forall\vec x.(A\Rightarrow B)$. Using All-E, we instantiate the bound variables $\vec x$ by fresh variables also named $\vec x$ (without loss of generality) and obtain $\Theta\vdash_\CL A\Rightarrow B$. Using Weak and Imp-E, we get $\Theta,A\vdash_\CL B$. Applying Theorem42, we derive $\Theta^T,A^T\vdash_\IL B^T$, that is $\Theta^T,A\vdash_\IL(B\Rightarrow R)\Rightarrow R$. By Lemma44, we derive $\Theta,A\vdash_\IL(B\Rightarrow R)\Rightarrow R$. Choosing $R:=B$, we get $\Theta,A\vdash_\IL B$. Using Imp-I and All-I, we derive $\Theta\vdash_\IL\forall\vec x.(A\Rightarrow B)$. We conclude that $\vdash_\IL\forall\vec x.(A\Rightarrow B)$ in $\mathcal T$.
\end{proof}
\section*{Exact author reference for finite translation cases}
[38] Thomas Traversi\'e. Monad Translations for Higher-Order Logic. FSCD2025, LIPIcs337, Article34, pp34:1--14. DOI10.4230/LIPIcs.FSCD.2025.34. The complete directly cited Lemma16, its proof, and all connected author Sections2--6 (definitions, finite rules, monad embedding and constructive translation proofs), together with the original bibliography, are retained below as support only; they are not additional counted targets.
